\documentclass[10pt]{article}
\usepackage[margin=0.95in]{geometry}
\usepackage{amsmath,amssymb}
\usepackage{booktabs}
\usepackage{xcolor}
\usepackage[colorlinks=true,linkcolor=blue,citecolor=blue]{hyperref}
\newcommand{\yes}{\textcolor{teal}{\textbf{supported}}}
\newcommand{\no}{\textcolor{red}{\textbf{refuted}}}
\newcommand{\partialv}{\textcolor{orange}{\textbf{partial}}}
\title{Do our scripted (MP) results follow Sokoli\'c et al., Dinh et al.,
and Rahaman et al.?\\
{\large A pre-specified test suite over three papers, with verdicts and
the cells that failed}}
\date{2026-07-25}
\begin{document}
\maketitle

\paragraph{Scope.} ``Scripted data'' means the MP collector only (SE(3)
min-jerk plan, feedforward/feedback tracking, reactive contact handoff);
the pointing-label dataset is retired. Six trained arms: $\{$L2, HT, HG,
MIP$\}$ at 200 demos and $\{$L2, MIP$\}$ at 2000 demos, all 300k steps,
one seed each. Success rate (SR) is the twofactor harness, AS$=8$, 60
seeds. All Jacobians below are of the \emph{deployed} policy map
(observation window $\to$ executed 8-step action chunk, through the full
sampler) unless stated otherwise.

\paragraph{The two theories.}
\textbf{(A) Sokoli\'c, Giryes, Sapiro, Rodrigues, \emph{Robust Large
Margin Deep Neural Networks} (arXiv:1605.08254)}: the generalization
error is controlled by the \emph{spectral norm of the network Jacobian in
the neighbourhood of the training samples} (divided by the classification
margin), with no dependence on width or depth; they also propose a
Jacobian-norm regularizer.
\textbf{(B) Dinh, Pascanu, Bengio, Bengio, \emph{Sharp Minima Can
Generalize For Deep Nets} (arXiv:1703.04933)}: for rectifier networks,
parameter-space flatness/sharpness measures can be changed arbitrarily by
function-preserving reparameterizations, so they cannot by themselves
explain generalization; any meaningful measure must be invariant.

\section{Theory A: the Jacobian spectral norm}

\subsection*{A1 --- cross-arm test}
Two independent measurements were made. Both use the deployed map, but
the first drew states and gaps from separate probes; the second
(\texttt{probe\_theory\_suite}) measures $\|J\|_2$, the train error and
the held-out error in \emph{one} pass with a single state protocol. We
report the second as primary.

\begin{center}\small
\begin{tabular}{lrrrrrr}
\toprule
arm & $\|J\|_2$ & $N_{\text{win}}$ & $\|J\|_2/\sqrt{N}$ & gap & held-out err & SR \\
\midrule
MIP-2k  & 3.86 & 74111 & 0.0142 & 0.00428 & 0.00464 & 100 \\
L2-2k   & 4.24 & 74111 & 0.0156 & 0.00433 & 0.00546 & 94 \\
HT-200  & 2.18 & 7431  & 0.0253 & 0.01169 & 0.01179 & 75 \\
MIP-200 & 2.38 & 7431  & 0.0276 & 0.01198 & 0.01211 & 94 \\
L2-200  & 2.99 & 7431  & 0.0347 & 0.01141 & 0.01192 & 63 \\
HG-200  & 3.08 & 7431  & 0.0357 & 0.01084 & 0.01107 & 87 \\
\bottomrule
\end{tabular}
\end{center}

$r(\|J\|_2/\sqrt N,\ \text{gap}) = +0.857$ and
$r(\|J\|_2/\sqrt N,\ \text{held-out error}) = +0.869$ --- but the
correlation is carried \emph{entirely by the data-scale contrast}: the
two 2k arms have small gaps and small $\|J\|_2/\sqrt N$, the four
200-demo arms have large ones. \textbf{Within the 200-demo group the
relation vanishes}: sorted by $\|J\|_2$ the gaps read
$2.18\!\to\!0.0117$, $2.38\!\to\!0.0120$, $2.99\!\to\!0.0114$,
$3.08\!\to\!0.0108$ --- i.e.\ flat to slightly \emph{decreasing}. The
apparent cross-arm success of Theory A is a two-group effect that
$\sqrt N$ already accounts for. \no{} as an explanation of the
objective axis; \partialv{} as a description of the data axis.

\emph{Two recorded corrections.} (i) A first pass using Frobenius norms
at rollout states gave $r=0.966$ and a 6/6 monotone ordering; it does not
survive the single-protocol measurement. (ii) In that first pass MIP-200's
gap ($0.0073$) was below L2-200's ($0.0095$); under the controlled
protocol the ordering \emph{reverses} ($0.0120$ vs $0.0114$). The
between-arm gap differences at 200 demos are within the state-sampling
noise, so no claim should rest on them.

\subsection*{A2 --- within-arm trajectory (no cross-arm confound)}
Tracking both quantities across snapshots of the \emph{same} arm removes
every cross-arm confound. The spectral norm falls by a factor of
$2.7$--$4.7$ over training in every arm; the gap does not follow it.

\begin{center}\small
\begin{tabular}{lrrrrr}
\toprule
snapshot & 20k & 60k & 120k & 180k & 300k \\
\midrule
L2-200 \ $\|J\|_2$ & 7.93 & 5.15 & 3.14 & 2.94 & 2.99 \\
L2-200 \ gap       & 0.0103 & 0.0127 & 0.0090 & 0.0101 & 0.0114 \\
MIP-200 $\|J\|_2$  & 7.38 & 2.51 & 2.40 & 2.37 & 2.38 \\
MIP-200 gap        & 0.0089 & 0.0092 & 0.0136 & 0.0125 & 0.0120 \\
L2-2k \ $\|J\|_2$  & 13.76 & 13.05 & 5.40 & 4.61 & 4.24 \\
L2-2k \ gap        & 0.0027 & 0.0038 & 0.0034 & 0.0035 & 0.0043 \\
\bottomrule
\end{tabular}
\end{center}

In L2-200 the norm falls $2.7\times$ while the gap rises slightly; in
MIP-200 it falls $3.1\times$ while the gap rises $35\%$; in L2-2k it
falls $3.2\times$ while the gap rises $60\%$. \no{} --- within a training
run the Jacobian spectral norm and the generalization gap move in
opposite directions.

\subsection*{A3 --- their own regularizer, correctly dosed}
Sokoli\'c et al.\ propose penalising the Jacobian norm. Our first attempt
($\lambda=10^{-3}$) destroyed the fit (SR 2/60, on-support error $+35\%$)
and is a void cell. Two correctly-dosed arms ($\lambda = 10^{-5},
10^{-4}$) are training; the prediction is a \emph{lower} $\|J\|_2$ and a
\emph{smaller} gap, with the fit intact.

\subsection*{A4 --- width independence}
The bound is claimed independent of width. A $2\times$-width L2 arm is
training; the prediction is that its gap follows its measured $\|J\|_2$,
not its parameter count.

\section{Theory B: reparameterization invariance}

The encoder is an MLP with LeakyReLU, which is positively homogeneous, so
scaling the units of a hidden layer by $\alpha$ and dividing the next
layer's incoming weights by $\alpha$ leaves the function \emph{exactly}
unchanged. Applied to L2-200 (measured function change: mean
$3.5\times10^{-5}$, on an action scale of $0.62$ --- float error):

\begin{center}\small
\begin{tabular}{lrrl}
\toprule
statistic & before & after & verdict \\
\midrule
encoder weight stable rank & [18.8, 21.0, 21.1] & [4.7, 7.9, 21.1] &
\textcolor{red}{gauge artifact} \\
embedding covariance PR     & 3.246 & 3.247 & invariant \\
encoder-Jacobian PR         & 4.609 & 4.609 & invariant \\
deployed-map PR             & 1.9460 & 1.9461 & invariant \\
deployed $\|J\|_2$          & 2.9809 & 2.9809 & invariant \\
\bottomrule
\end{tabular}
\end{center}

\paragraph{Consequence for our own evidence.} We had cited the encoder
weight stable rank as data-free evidence that the L2/MIP split ``lives in
the weights''. That claim is \no{} and is withdrawn: the same function
admits stable ranks of $18.8$ and $4.7$.

\paragraph{What survives, and why.} The embedding PR survives this
symmetry because the conditioning path applies a pointwise GELU
\emph{directly} to the embedding
(\texttt{cond\_encoder = Sequential(GELU, Linear)}), which pins the
embedding's coordinate system --- it is not a free gauge in this
architecture. It is, however, only weakly protected: rescaling the
embedding's output units and compensating in the (GELU-preceded) consumers
moved emb\_PR from $3.25$ to $4.84$ ($+49\%$) while changing the function
by $\sim1\%$. We therefore treat emb\_PR as admissible but flagged, and
treat the deployed-map quantities --- which are function-space objects and
invariant by construction --- as the primary evidence.

\subsection*{B2 --- $\epsilon$-sharpness}
Dinh's own quantity
$S(\theta)=\max_{\|\delta\|\le\epsilon\|\theta\|}
(L(\theta+\delta)-L(\theta))/(1+L(\theta))$, $\epsilon=5\times10^{-4}$,
16 random directions, measured before and after the \emph{exact}
function-preserving rescaling:

\begin{center}\small
\begin{tabular}{lrrr}
\toprule
arm & $S(\theta)$ & after rescaling & ratio \\
\midrule
L2-200  & $1.93\times10^{-3}$ & $5.19\times10^{-4}$ & $0.27\times$ \\
L2-2k   & $3.95\times10^{-3}$ & $1.35\times10^{-3}$ & $0.34\times$ \\
HT-200  & $1.50\times10^{-2}$ & $2.55\times10^{-2}$ & $1.70\times$ \\
HG-200  & $1.19\times10^{-1}$ & $3.26\times10^{-2}$ & $0.27\times$ \\
MIP-200 & $6.25\times10^{-2}$ & $8.09\times10^{-2}$ & $1.29\times$ \\
MIP-2k  & $6.95\times10^{-2}$ & $0$                 & $0\times$ \\
\bottomrule
\end{tabular}
\end{center}

Sharpness moves by factors of $0.27$ to $1.7$ (and to exactly zero for
MIP-2k) under a transformation that leaves every prediction unchanged.
\yes{} --- Dinh's construction applies to our networks directly.

Two further observations. First, the deployed $\|J\|_2$ and deployed PR
are invariant to \emph{machine precision} under both the rescaling and an
exact hidden-unit permutation (e.g.\ L2-200: $2.9903 \to 2.9903 \to
2.9903$), confirming the function-space/parameter-space split.
Second, and directly in the spirit of the paper's title: MIP-200 sits in a
minimum $32\times$ \emph{sharper} than L2-200's ($6.3\times10^{-2}$ vs
$1.9\times10^{-3}$) while generalizing at least as well and succeeding
$50\%$ more often. Sharp minima generalize here too. (Caveat: the loss
functions differ between arms, so cross-arm sharpness comparisons are
suggestive, not rigorous; the within-arm reparameterization results are
the rigorous part.)

\section{Theory C: spectral bias (Rahaman et al., arXiv:1806.08734)}

\emph{Claim.} ReLU networks learn low frequencies first and express high
frequencies only with finely tuned parameters, so high-frequency content
is acquired late and is fragile.

\emph{Test.} A demonstration trajectory is a dense 1-D path through
observation space with ground-truth actions at every point, so the
frequency content of a policy's output along that path can be compared
directly with the target's own spectrum. Over 12 trajectories per arm we
report the captured power ratio (pred/GT) in three bands
($f<0.05$, $0.05\le f<0.2$, $f\ge0.2$ cycles/step), the error power by
band, and the high-band power lost under a small parameter perturbation.
The target itself is strongly low-frequency: its high/low power ratio is
$0.0102$.

\begin{center}\small
\begin{tabular}{lrrrrr}
\toprule
 & \multicolumn{3}{c}{captured high-band power ratio} & err.\ frac.\ & frag.\ \\
arm & 20k & 60k & 300k & high @300k & high \\
\midrule
L2-200  & 0.901 & 0.913 & 0.996 & 0.587 & $+0.001$ \\
MIP-200 & \textbf{0.978} & \textbf{0.993} & 0.988 & 0.634 & $-0.004$ \\
HT-200  & 0.919 & 0.958 & 0.988 & 0.624 & $-0.009$ \\
HG-200  & 0.752 & 0.891 & 0.984 & 0.728 & $+0.001$ \\
L2-2k   & 0.723 & 0.860 & 0.993 & 0.383 & $+0.001$ \\
MIP-2k  & \textbf{1.012} & 0.964 & 1.000 & 0.352 & $+0.000$ \\
\bottomrule
\end{tabular}
\end{center}

\paragraph{Learning order: \yes{}.} High-band capture rises monotonically
with training in four of six arms (HG $0.75\!\to\!0.98$, L2-2k
$0.72\!\to\!0.99$), i.e.\ low frequencies are fitted first --- exactly the
predicted order. The residual error is high-frequency dominated in every
arm ($35$--$73\%$ of error power above $f=0.2$), again as predicted.

\paragraph{A new objective difference.} \textbf{MIP acquires the
high-frequency content first}: at 20k steps it is already at $0.978$
(200 demos) and $1.012$ (2k) while HG is at $0.752$ and L2-2k at
$0.723$. This lines up with two independent measurements: MIP rebuilds
representation rank from $\sim$1000 steps (§B), and MIP's success rate
peaks at $20$--$60$k. Fast acquisition of fine structure is a genuine,
measurable property of the objective.

\paragraph{Endpoint deficiency: \no{}.} By 300k \emph{every} arm captures
$98$--$100\%$ of the target's high-band power. The SR spread of $63$--$100$
therefore cannot be attributed to missing high-frequency content along the
demonstration manifold.

\paragraph{Fragility: not observed.} A relative parameter perturbation of
$2\times10^{-3}$ changes high-band power by less than $1\%$ in every arm
($|{\rm frag}| \le 0.009$), so the ``finely tuned parameters'' claim does
not manifest at this scale in our networks.

\section{What the theories do and do not explain}

\paragraph{Explained.} Theory C's learning-order claim is confirmed, and
it yields a new and useful fact: MIP reaches the target's high-frequency
content several times faster than the other objectives, which coincides
with its early SR peak. Theory B is confirmed outright and gives us a
measurement filter we now apply: \emph{function-space quantities only}.
It also retires two things we had been quoting --- encoder weight stable
rank (a gauge artifact) and any prospective ``MIP finds flatter minima''
story (MIP's minima are $32\times$ sharper). Theory A survives only as a
statement about the \emph{data} axis: $\|J\|_2/\sqrt N$ separates the 2k
arms from the 200-demo arms in the right direction.

\paragraph{Not explained.} (i) The objective axis: at matched data the
four objectives span $63$--$94$ SR with $\|J\|_2$ spanning only
$2.18$--$3.08$ in the \emph{wrong} order (HT lowest norm, second-worst
SR; HG highest norm, third-best SR) and with held-out errors that agree
to within sampling noise ($0.0111$--$0.0121$). Neither theory addresses
why these four functions, which are equally accurate off their training
sets, behave so differently in closed loop. (ii) The \emph{formation} asymmetry: from an
identical initialization (emb\_PR $4.24$, encoder-Jacobian PR $35.7$),
every objective collapses to emb\_PR $\approx2.6$--$3.1$ within
$25$--$100$ steps, after which HG/HT rebuild from $\sim$400 steps, MIP
from $\sim$1000, and L2 barely rebuilds at all
($2.65\!\to\!3.69$ by 20k versus HG's $2.61\!\to\!6.54$). Neither theory
addresses which objectives rebuild.

\paragraph{Interventions attempted (honest accounting).} Five arms
designed to move a candidate property have been damaged by dose or by
train/test mismatch and are void cells, not evidence: obs-noise
($\sigma=0.05$, SR 12), jacreg ($\lambda=10^{-3}$, SR 2), condreg v1
($\lambda=1$, SR 0), featdrop $p=0.05$ (SR 17) and $p=0.15$ (SR 2). The
featdrop pair is instructive nonetheless: it \emph{did} raise emb\_PR
($3.29\to4.23$) and it \emph{did} shrink the gap ($0.0095\to0.0063$),
while SR collapsed --- because the level of the held-out error rose
($0.0100\to0.0136$). SR follows the error level, not the gap.
Correctly-dosed replacements (condreg v2 at $3\times10^{-3}$ and
$3\times10^{-4}$, jacreg at $10^{-5}$ and $10^{-4}$) are training.

\section{The experiments run to try to establish the account}

Every measurement below was undertaken specifically to give the three
theories the best chance of explaining MIP/flow's advantage.

\subsection*{E1 --- Sokoli\'c where the margin is tight (insertion phase)}
The bound is margin/$\|J\|_2$; our margin is smallest at insertion, so that
is where the theory should bite. A first pass ($n=40$ states) looked
decisive --- MIP's insertion-phase gain appeared $3\times$ below L2's at
both data scales ($5.6$ vs $16.2$; $10.3$ vs $35.6$). \textbf{It does not
replicate at $n=120$}: the means are outlier-dominated and unstable
(MIP-200 now reads $56.3$, \emph{above} L2-200's $26.3$). Robust
statistics give

\begin{center}\small
\begin{tabular}{lrrr}
\toprule
arm & gain median & gain p90 & SR \\
\midrule
MIP-200 & 2.40 & \textbf{18.6} & 94 \\
HT-200  & 2.71 & 54.4 & 75 \\
HG-200  & 3.11 & 67.4 & 87 \\
L2-200  & 3.11 & 49.4 & 63 \\
L2-2k   & 4.10 & 74.2 & 94 \\
MIP-2k  & 4.47 & \textbf{23.7} & 100 \\
\bottomrule
\end{tabular}
\end{center}

$r=+0.18$ (median) and $-0.02$ (p90) against SR. The one real signal is
that \emph{both MIP arms have far fewer high-gain insertion states}
(p90 $18.6$/$23.7$ versus $49$--$74$ for every other arm) --- an objective
signature, but one that does not order the remaining four arms. \no{} as
stated; \partialv{} as a MIP-vs-rest contrast.

\subsection*{E2 --- Rahaman off the manifold (pre-registered falsifier)}
On-manifold every arm is spectrally complete, so the pre-registered test
was whether fine structure survives \emph{off} it. Along captured
off-support rollout states ($d>10$, 600 states) versus on-support ones:

\begin{center}\small
\begin{tabular}{lrrrr}
\toprule
arm & high-band frac.\ on & off & output power off/on & SR \\
\midrule
MIP-200 & 0.607 & 0.559 & \textbf{0.29} & 94 \\
MIP-2k  & 0.610 & 0.597 & \textbf{0.30} & 100 \\
L2-2k   & 0.610 & 0.605 & 0.49 & 94 \\
HT-200  & 0.606 & 0.586 & 0.57 & 75 \\
HG-200  & 0.608 & 0.582 & 0.71 & 87 \\
L2-200  & 0.606 & 0.590 & \textbf{1.03} & 63 \\
\bottomrule
\end{tabular}
\end{center}

The spectral \emph{shape} is preserved by every arm off-manifold
($0.56$--$0.60$ high-band fraction, essentially unchanged). What differs is
the \emph{total output power}: the MIP arms attenuate to $29$--$30\%$ of
their on-manifold energy, L2-2k to $49\%$, and L2-200 not at all
($103\%$). This orders SR $5/6$ with $r=-0.844$ --- the best of any
quantity in this report --- but it is a statement about response
\emph{magnitude} off the data, i.e.\ a re-measurement of the off-support
behaviour we already had ($|a|$ far-field $0.40$ for MIP vs $0.79$ for
L2), not a spectral-bias effect.

\subsection*{E3 --- causal tests of the in-support account}
The incoherent (within-chunk) component of the held-out error orders SR
$6/6$ --- the only on-manifold quantity that separates HT from MIP. Two
manipulations tested it causally. Injecting incoherent error into MIP-200
at doses spanning the entire between-arm range ($0.05$--$0.40$mm, both
i.i.d.\ and the faithful state-correlated zero-mean-per-chunk version)
left SR unchanged: $52/60 \to 51,54,50,51,52,57$. Removing it from L2-200
by within-chunk low-pass destroyed the policy (the filter attenuates the
coherent component too) and is void. \no{} --- the correlation does not
survive intervention; the PD/impedance layer and arm inertia evidently
absorb incoherent action error at these magnitudes.

\subsection*{E4 --- ablations}
Batch size $256$ vs $1024$: SR $41/60$, $38/60$ vs $40/60$, $38/60$; in-support
error identical ($1.054$ vs $1.046$mm). The gradient-noise knob does
nothing, weakening every noise-driven account. Coverage: $200 \to 38/60$,
$800 \to 55/60$, $2000 \to 56/60$, with in-support error falling
$1.046 \to 0.754 \to 0.449$mm --- on the \emph{data} axis precision and SR
move together, which is what an on-manifold theory predicts and is the
one place the theories work.

\section{Verdict}

\begin{center}\small
\begin{tabular}{llp{7.4cm}}
\toprule
theory & verdict & what it does and does not carry \\
\midrule
Dinh & \yes{} & Applies to our networks exactly; retires weight stable
rank and any flat-minima story (MIP's minima are $32\times$ sharper);
supplies the function-space measurement filter. Methodological, not
explanatory. \\
Sokoli\'c & \partialv{} & Describes the \emph{data} axis (2k vs 200).
Fails within a training run (norm falls $2.7$--$4.7\times$ while the gap
is flat or rises), fails within matched data, and its margin-weighted
form fails at adequate $n$. \\
Rahaman & \partialv{} & Learning order confirmed, and MIP acquires
high-frequency content $\sim3\times$ earlier --- a real objective
difference. But every arm is spectrally complete by 300k, so it cannot
explain the endpoint spread. \\
\bottomrule
\end{tabular}
\end{center}

\paragraph{The honest conclusion.} The three papers are theories of
on-manifold supervised learning, and on that ground our six arms differ
in the ways the papers describe --- but those differences do not survive
intervention, while the differences that track SR (off-support response
magnitude, $r=-0.84$; excursion escalation, $0.44$ vs $0.12$; the
restoring-versus-escaping annulus field) are all off-manifold quantities
that none of the three papers addresses. We therefore cannot claim these
papers explain MIP/flow's benefit. What we can claim is narrower and
defensible: \emph{Dinh tells us which of our measurements are admissible;
Sokoli\'c accounts for why more data helps; Rahaman accounts for why MIP
gets there first; and the residual advantage --- the part that separates
the objectives at fixed data --- is an off-manifold property that remains
unexplained by all three.}

\section{Pre-registered falsifiers}
\begin{enumerate}\itemsep2pt
\item If the within-arm trajectory (A2) shows the gap and $\|J\|_2$ moving
in opposite directions over training, Theory A does not describe our
training dynamics, only our cross-arm variation.
\item If correctly-dosed jacreg (A3) lowers $\|J\|_2$ with the fit intact
and the gap does \emph{not} shrink, Theory A's mechanism is refuted
constructively in our setting.
\item If the $2\times$-width arm (A4) has the same $\|J\|_2$ and a
different gap, the width-independence claim fails here.
\item If $\epsilon$-sharpness is unchanged by the exact rescaling (B2),
our architecture escapes Dinh's construction and sharpness would become
admissible again.
\item If a frequency-resolved measurement \emph{off} the demonstration
manifold (the regime where the arms actually differ) shows the spectral
deficiency that the on-manifold measurement does not, Theory C would be
reinstated as an explanation of the objective axis rather than only of
the learning order.
\end{enumerate}

\end{document}
